\leavevmode{\parindent=1em\indent} Trong bối cảnh công nghệ thông tin ngày càng phát triển mạnh mẽ, việc ứng dụng công nghệ vào giải quyết các bài toán thực tiễn đóng vai trò quan trọng. Trong đó, nhiều văn hoá truyền thống của dân tộc Việt Nam được ``số hoá'', tạo cơ hội cho những người trẻ - những người sống trong thời đại công nghệ được tiếp cận với truyền thống văn hoá nước nhà theo một cách mới mẻ và thú vị hơn.

\leavevmode{\parindent=1em\indent} Xuất phát từ ý tưởng số hóa và tổ chức thông tin về các mối quan hệ huyết thống trong gia đình, nhóm chúng em lựa chọn thực hiện đề tài \textbf{``Chương trình Quản lý gia phả''}. Mục tiêu của chương trình là xây dựng một hệ thống cho phép lưu trữ, quản lý và truy vấn thông tin các thành viên trong gia đình dưới dạng cấu trúc cây, từ đó hỗ trợ thực hiện các thao tác như thêm thành viên, tìm kiếm quan hệ họ hàng, duyệt cây gia phả và hiển thị thông tin một cách có hệ thống. Thông qua việc xây dựng chương trình, nhóm có cơ hội vận dụng các kiến thức cơ bản về cấu trúc dữ liệu, thuật toán và lập trình để giải quyết bài toán quản lý thông tin có cấu trúc phân cấp - đúng như tinh thần của phương pháp \textit{học tập qua dự án}.

\leavevmode{\parindent=1em\indent} Trong quá trình thực hiện dự án, nhóm đã nhận được sự hướng dẫn của TS. Nguyễn Văn Hiệu, giảng viên phụ trách học phần. Nhóm xin chân thành cảm ơn thầy đã định hướng và góp ý để nhóm có thể hoàn thành dự án này.

\leavevmode{\parindent=1em\indent} Mặc dù rất mong muốn thực hiện một chương trình hoàn chỉnh cả về nền tảng lý thuyết và phương pháp thực hiện, cũng như mong muốn có thêm phản hồi từ người dùng về chất lượng chương trình, tuy nhiên vì thời gian và trình độ còn hạn chế, vì vậy chương trình và báo cáo chắc chắn sẽ còn tồn tại những thiếu sót nhất định, kính mong quý thầy cô đóng góp ý kiến cho chương trình và báo cáo để nhóm rút kinh nghiệm cho những dự án PBL tiếp theo trong quá trình học tập tại trường.

\leavevmode{\parindent=1em\indent} Xin trân trọng cảm ơn.
